\chapter{PlusCal / TLA+ Quick Reference}
\label{chap:pluscal-quick-ref}

This appendix is the working reference behind every \emph{[verification]}
exercise in the book.  It assumes Chapter~6 (\S\ref{sec:pluscal-example},
\S\ref{sec:verify-workflow}); here you find the syntax, the TLC commands,
and the failure modes without re-reading the chapter.  TLA+ is the language
of state machines; PlusCal is the pseudo-code layer that compiles to it.

\section{The Language at a Glance}

A TLA+ specification names \emph{variables}, sets their \emph{initial
values}, and lists \emph{actions} that transform one state into the next.
The temporal formula \texttt{Spec == Init \textbackslash/\ [][Next]\_vars}
says: start at \texttt{Init}, then take \texttt{Next} (or stutter)
forever.  TLC then enumerates the reachable states.

\begin{lstlisting}[language={}, caption={A minimal TLA+ state machine.}]
---- MODULE M ----
EXTENDS Naturals, TLC
VARIABLES x
Init == x = 0
Inc == x' = x + 1
Next == Inc
Spec == Init /\ [][Next]_x
====
\end{lstlisting}

\emph{Constants} are declared with \texttt{CONSTANTS} and fixed in the
configuration file.  An \emph{invariant} is a state predicate that must hold
in every reachable state; a \emph{property} (temporal) asserts that a good
state is eventually reached under stated fairness assumptions.  The table in
\S\ref{sec:verify-workflow} maps every runtime backend in this book to its
model file, key invariant, and TLC outcome.

\section{Writing a Configuration File}

The configuration is part of the proof artifact.  It names the constants,
the invariants, and any temporal property, and it keeps the state space
finite.  The canonical form (matching \S\ref{sec:pluscal-example}) is:

\begin{lstlisting}[language={}, caption={A TLC configuration file.}]
SPECIFICATION Spec

CONSTANTS
  Accounts = 1..2
  Floor = 0

INVARIANT
  TypeOK
  MoneyConserved
  FloorsHold

CHECK_DEADLOCK TRUE
\end{lstlisting}

\emph{CONSTRAINT} entries further bound the state space (e.g. a maximum
clock value) when an unbounded counter makes TLC explore forever; see the
\texttt{TLCBound} pattern used by the SPHT and TiKV models in
\texttt{docs/proofs/}.  \texttt{CHECK\_DEADLOCK TRUE} turns deadlock into an
error, which is how \emph{[verification]} exercises about eager locking
find their deadlock.

\section{A TM Model Skeleton}

Every \emph{[verification]} exercise in this book is a small variation of
the same skeleton: two accounts, a global clock or lock, and the two
invariants \texttt{MoneyConserved} and \texttt{FloorsHold}.  The PlusCal
form (see \S\ref{sec:pluscal-example} for the full running model):

\begin{lstlisting}[language={}, caption={The shared shape of every TM model in the book.}]
--algorithm Bank
variables
  clock = 0, lock = FALSE,
  bal = [a \in 1..2 |-> 10];

process Tx \in 1..2
variables snap = 0, r1 = 0, r2 = 0;
begin
  BeginTx:  snap := clock;
  Read:     r1 := bal[1]; r2 := bal[2];
  Validate: if clock /= snap then abort or re-read; end if;
  Commit:   await lock = FALSE; lock := TRUE;
            bal[1] := bal[1] - 1; bal[2] := bal[2] + 1;
            clock := clock + 1; lock := FALSE;
end process;
end algorithm;
====
\end{lstlisting}

The three layers of state from \S\ref{sec:pluscal-example} --- application
state (\texttt{bal}), transaction state (\texttt{snap}, \texttt{r1},
\texttt{r2}), and metadata (\texttt{lock}, \texttt{clock}) --- are always
kept separate.  An \emph{opacity} model adds a read set and a
commit-time recheck; \S\ref{sec:norec-example} and
\S\ref{sec:pluscal-bank-floor-example} show both validations.

\section{Running TLC}

The companion repository's \texttt{docs/proofs/Makefile} wraps TLC.  After
downloading the checker once (\texttt{make download-jar}), a single backend
is checked with:

\begin{verbatim}
java -cp /tmp/tla2tools.jar tlc2.TLC -deadlock \
     Bank.tla -config Bank.cfg
\end{verbatim}

Useful flags: \texttt{-workers N} for parallel checking;
\texttt{-coverage} to report which actions actually fired (a warning sign
when an action never fires); \texttt{-maxSetSize N} to bound huge sets.
Inside the repository, \texttt{make -C docs/proofs check} runs every model
in one pass, and \texttt{make -C docs/proofs verify-liveness} adds the
fairness-based properties.

\section{Pitfalls}

\begin{itemize}
    \item \textbf{Unbounded state spaces.}  A counter that only increments
          (the commit clock, retry counters, sequence numbers) makes TLC run
          forever.  Bound it with a \texttt{CONSTRAINT} or a small
          \texttt{CONSTANTS} domain; the book's configs use finite
          \texttt{TLCBound} constraints for exactly this reason.
    \item \textbf{Record-valued state.}  Function-valued variables such as
          \texttt{bal = [a $\in$ 1..2 |-> 10]} fingerprint fine in TLC, but a
          model with per-process records that is translated from PlusCal can
          blow up because each assignment creates a new record value.
          Prefer named scalars and small sets.
    \item \textbf{Missing labels.}  PlusCal requires a label before every
          statement that touches a shared variable or that other processes
          must await; the translator errors with ``missing label'' when a
          loop or a multi-statement block lacks one.
    \item \textbf{Type errors hide real errors.}  \texttt{TypeOK} (a
          predicate over the types of every variable) is the first
          invariant to write; without it, a modelling mistake surfaces as a
          confusing \texttt{Invariant \ldots\ is violated} on an unrelated
          formula.
    \item \textbf{Stuttering.}  A variable missing from \texttt{UNCHANGED}
          in an action changes the set of legal transitions.  When an
          exercise asks you to split an action into \textsc{Debit} and
          \textsc{Credit} steps, the point is precisely to observe the
          intermediate states that a single atomic action hides.
\end{itemize}

\section{Reading a Counterexample}

A TLC counterexample is a schedule, and it is read forward.  Start at the
initial state and step through each state transition:

\begin{enumerate}
    \item Identify the process that moved and the label that fired.
    \item Identify the shared variable that changed (the balance, the
          clock, the lock, the read set).
    \item Ask which earlier read became stale at that point --- the
          \emph{[verification]} exercises in this book are almost always
          about a stale read that validation failed to reject.
    \item Replay the same schedule mentally on the runtime code: the model
          step corresponds to a source-level read/write pair (see
          \S\ref{sec:tla-bank-conservation} for the concrete mapping).
\end{enumerate}

If the injected-bug exercises (Chapter~6, ``Bug-Injection'') do not
produce a counterexample, the mutation was either irrelevant or the
invariant is too weak --- check both before concluding the model is right.